%% %%%%%%%%%%%%%%%%%%%%%%%%%%%%%%%%%%%%%%%%%%%%%%%%%
%% Template for working papers
%% Tommaso Di Francesco
%% University of Amsterdam
%% %%%%%%%%%%%%%%%%%%%%%%%%%%%%%%%%%%%%%%%%%%%%%%%%%
%% Version 1.0 // February 2024
%% 
\documentclass[11pt]{article}
\usepackage{style}


%% ===============================================
%% Setting the line spacing (3 options: only pick one)
% \doublespacing
% \singlespacing
\onehalfspacing
%% ===============================================

\setlength{\droptitle}{-5em} %% Don't touch

% %%%%%%%%%%%%%%%%%%%%%%%%%%%%%%%%%%%%%%%%%%%%%%%%%%%%%%%%%%
% SET THE TITLE
% %%%%%%%%%%%%%%%%%%%%%%%%%%%%%%%%%%%%%%%%%%%%%%%%%%%%%%%%%%
% TITLE:
\title{Lasso Trading}
\author{}
% DATE:
\date{\today}

% %%%%%%%%%%%%%%%%%%%%%%%%%%%%%%%%%%%%%%%%%%%%%%%%%%%%%%%%%%
% %%%%%%%%%%%%%%%%%%%%%%%%%%%%%%%%%%%%%%%%%%%%%%%%%%%%%%%%%%
\begin{document}
% %%%%%%%%%%%%%%%%%%%%%%%%%%%%%%%%%%%%%%%%%%%%%%%%%%%%%%%%%%
% %%%%%%%%%%%%%%%%%%%%%%%%%%%%%%%%%%%%%%%%%%%%%%%%%%%%%%%%%%

{\setstretch{.8}
\maketitle
\centering
}

% \section{Introduction}

% The volume and granularity of data available worldwide is growing at an unprecedented pace.
% In 2024 alone, the world generated approximately 402 million terabytes of data per day, up from 43 million just a decade earlier \citep{DataSphere2025}. 
% Financial markets are no exception to this trend.
% Whereas investors once had access to only a handful of data series, today they can draw on a vast array of signals, ranging from traditional balance-sheet information to textual data, satellite imagery, and web traffic.
% For most of these new variables, there is no clear theoretical prior about their relationship, if any, to future stock returns. 
% Nevertheless, such data are in high demand with investors. 
% The reason is that investors increasingly rely on machine learning and artificial intelligence (AI) to find useful patterns in large datasets.\footnote{The Invesco Global Systematic Investing Study 2024 reports that more than half of surveyed investors incorporate AI, using tools ranging from pattern recognition to portfolio optimization \citep{Invesco2024IGSIS}.}
% These methods require no theorytical guidance about which variables should matter; instead, they let the data speak for themselves.

% This behavior is not well captured by existing asset pricing models with learning. 
% In these models, agents typically are assumed to employ self-referential or correctly specified forecasting rules, or to work with severely under-parameterized perceived laws of motion. Agents track a small number of parameters and update them recursively using variants of least squares or constant-gain learning. 
% These frameworks are analytically convenient, but they do not describe the behabior of modern invstors investors who consider thousands of potential signals to forecast returns and use flexible statistical procedures to discover such relationships without an explicit structural model in mind.

% We develop a model that embeds this form of statistical learning into a standard asset pricing environment. 
% Under rational expectations, returns in our economy are not predictable. 
% Agents, however, suspect that returns may be conditionally predictable using a large set of potential predictors. 
% To exploit this perceived opportunity, they bet on sparsity and apply a machine learning method, least absolute shrinkage and selection operator (LASSO), to select a sparse subset of variables from the high-dimensional predictor space and to construct forecasting rules \citep{Tibshirani1996}.

% Our main result is that the very act of searching for predictability with such statistical learning methods can generate short-lived pockets of return predictability. 
% In equilibrium, these pockets are sparse and evolve over time, echoing recent empirical evidence on time-varying short-term predictability in financial markets \citep{Chinco2019,Farmer2023}. 
% In essence, the process of learning in high-dimensional data both hunts for and endogenously creates the temporary predictability it seeks to exploit, providing a new perspective on how modern data-driven investment practices interact with asset prices.

% We provide what is, to our knowledge\footnote{\cite{Georges2021} study by means of simulation various types of regularization methods, including LASSO
% in a model in which forecasting rules are auto-regressive including higher-rder polynomials of past returns to allow for non-linearities.
% Compared to them we provide analytical results and focus on the use of LASSO in linear regression on high-dimensional data.}, the first theoretical characterization of LASSO-based learning in an equilibrium asset-pricing environment. 
% This allows us to link a literature that has relied on low-dimensional, stylized learning rules with the empirical tradition in which researchers and practitioners routinely estimate high-dimensional predictive models.
% Endowing agents with a realistic learning rule is not only conceptually appealing; it also enables a tight connection between theory and data. 
% On the theoretical side, we show that LASSO learning generates pockets of predictability even when returns are fundamentally unpredictable. 
% We analytically characterize the probability with which such pockets arise, providing a precise mechanism that rationalizes the sparse and episodic nature of return predictability documented in recent literature. 
% This framework delivers strong predictions regarding when and how predictability emerges in finite samples, which we test using a wide set of assets and a rich vocabulary of potential predictors constructed from textual “topics” \citep{Bybee2021}.

% \textit{The Model} We study a standard Lucas-type exchange economy with a single risky asset that pays stochastic dividends and a representative consumption process. 
% Under rational expectations, the equilibrium price–dividend ratio is constant and returns are unpredictable. 
% We extend this framework along the lines of \cite{Adam2016} by assuming that investors fully understand the dividend and consumption processes but hold a misspecified view of the process that generates asset returns. 
% Agents believe that log returns may be predictable from observable signals. 
% At time \textit{t}, they use recent histories of returns and signals to estimate a perceived law of motion (PLM) in which log returns at \textit{t+1} depend linearly on a subset of available predictors. 
% Because the number of candidate signals is large relative to the amount of data available, agents rely on regularization to select which predictors to include. 
% We model this learning behavior using the least absolute shrinkage and selection operator (LASSO). 
% Agents’ beliefs feed back into prices through the actual law of motion (ALM) implied by equilibrium asset pricing. 
% Importantly, although agents employ linear forecasting rules, the ALM for returns is nonlinear. As a result, agents never learn the true data-generating process. 
% Instead, learning converges to a weaker notion of equilibrium based on statistical self-consistency: the moments implied by agents’ forecasting regressions coincide with those generated by equilibrium prices. 
% We formalize this concept as a cross-moment consistent equilibrium. 
% In such an equilibrium, agents’ estimated forecasting coefficients are, on average, those implied by the joint distribution of returns and predictors in equilibrium, even though returns remain fundamentally unpredictable. 
% The central mechanism of the model is simple. In finite samples, regularized learning occasionally selects predictors whose estimated coefficients are large enough to survive the LASSO penalty. 
% When this occurs, agents temporarily perceive return predictability and adjust their trading behavior accordingly. This trading feeds back into prices, generating short-lived but genuine return predictability. 
% As new data arrive and older observations are discarded or discounted, estimated coefficients revert toward zero, perceived predictability disappears, and prices again behave as if returns were unpredictable. 

% \textit{Results} Return predictability in our model emerges intermittently, is sparse across predictors, and short-lived. The model thus matches the empirical patterns documented in \cite{Chinco2019,Farmer2023}. 

% \textit{Related Literature} 
% Our model relates to the large literature on asset pricing with learning. 
% The central premise of this work is that investors have incomplete knowledge of the data-generating process governing market outcomes and therefore infer it from observed data. 
% Learning is often introduced to account for empirical features, such as excess volatility, return predictability, and time variation in price–dividend ratios, that are hard to replicate with fully informed agents. 
% Models differ in what agents know about the economic structure, the learning technology they use, and the objects they attempt to learn.

% Early contributions assume fully rational agents who estimate unknown parameters of a correctly specified model using least squares \citep{Marcet1989}. 
% In this setting, learning converges to the rational expectations equilibrium (REE) as agents observe more data.
% Even in this low-dimensional setting, learning can generate excess volatility and return predictability \citep{Timmermann1993,Timmermann1996}. 
% Convergence to the REE is no longer guaranteed when agents are boundedly rational; with finite-memory and finite-sample learning, beliefs may fail to converge \citep{Honkapohja2003}. 
% The implications become richer when agents face model uncertainty in addition to parameter uncertainty. 
% \cite{barberis1998model} study investors who switch between competing forecasting rules, while \cite{hong2007simple,Branch2010} analyze settings in which agents estimate univariate perceived laws of motion despite a multivariate true model.

% More recently, \cite{adam2011internal} introduced the concept of internal rationality to the learning literature.
% Internally rational agents optimize their decisions given their subjective beliefs, even if these beliefs are misspecified.
% Building on this insight, \cite{Adam2016} embed internally rational agents in an otherwise standard asset-pricing model.
% In their model, agents know the dividend and consumption processes but have to estimate the mapping from fundamentals to prices from observed data.
% This generates a feedback loop between beliefs and prices that helps match key asset-pricing moments. 
% Our model builds on \cite{Adam2016} because their setup features a very general learning mechanism that generates feedback between beliefs and prices. 
% This allows us to focus on the implications of replacing their general learning rule with a microfounded LASSO-based learning mechanism.

% Our paper also relates to research on how big data and improved statistical tools affect asset prices. 
% \cite{begenau2018big} show that higher data-processing capacity improves investor forecasts and changes equilibrium prices and the cost of capital, while \cite{farboodi2020long} model investors who optimally choose whether to aquire data about fundamentals or about the behavior of other investors.

% Empirical evidence indicates that return predictability is time-varying. 
% Variables such as consumption growth and valuation ratios forecast returns in some periods but not in others \citep{AngBekaert2007_RFS,GoyalWelch2003_MS,LettauLudvigson2001_JF}. 
% A complementary strand of research emphasizes sparsity: among the many predictors proposed in the literature, only a small and time-varying subset appears informative at any point in time \citep{kozak2020shrinking,freyberger2020dissecting,Feng2020}. 
% \cite{Mclean2016} show that many return anomalies attenuate after publication, consistent with investors actively searching for and trading on these signals.
% The ever-increasing number of published return predictors, sometimes coined as "factor zoo", recently raised concerns about data mining which motivated the use of machine learning in return prediction \citep{harvey2016and}. 
% \cite{gu2020empirical} show that machine-learning methods outperform traditional linear models by exploiting sparse signals. 
% Most closely related to our learning mechanism, \cite{Chinco2019} document that LASSO uncovers short-lived pockets of high-frequency predictability concentrated in a small, shifting set of predictors. 

% Finally, our work connects to initial attempts to model market outcomes when agents use statistical learning. 
% \cite{Georges2021} study adaptive model selection and regularization in an agent-based market and show that regularization mitigates but does not eliminate endogenous instability. 
% Relative to this simulation-based evidence, our contribution is to provide an analytical characterization of LASSO-based learning in an equilibrium asset-pricing environment and to derive sharp conditions under which sparse, temporary predictability arises endogenously even when returns are fundamentally unpredictable.

% The rest of the paper is structured as follows. Section 2 presents the model. 
% Section 3 analyzes the asset-pricing equilibrium under different learning rules. 
% Section 4 estimates the model on the cross-section of US stocks and a high-dimensional set of predictors derived from textual data.
% Section 5 concludes.
% =====================================================



\section{Introduction}

The volume and granularity of data available worldwide is growing at an unprecedented pace. In 2024 alone, the world generated approximately 402 million terabytes of data per day, up from 43 million just a decade earlier \citep{DataSphere2025}. 
Financial markets are no exception to this trend. 
Whereas investors once had access to only a handful of data series, today they draw on a vast array of signals, ranging from traditional balance-sheet information to textual data, satellite imagery, and web traffic.
For most of these new variables, there is little theoretical guidance about whether they should forecast stock returns. 
Nevertheless, demand for alternative data is strong, and investors increasingly use machine learning and artificial intelligence (AI) to search for predictive patterns in large datasets \citep{Invesco2024IGSIS}. 
These methods require no theoretical guidance about which variables should matter; instead, they let the data speak for themselves.

This behavior is not well captured by existing asset pricing models with learning. 
In these models, agents typically are assumed to employ self-referential or correctly specified forecasting rules, or to work with severely under-parameterized perceived laws of motion. 
Agents track a small number of parameters and update them recursively using variants of least squares or constant-gain learning.
While analytically convenient, these frameworks do not describe the behavior of modern investors who consider thousands of potential signals to forecast returns and use flexible statistical procedures to discover relationships without an explicit structural model in mind.

We develop a model that embeds this form of statistical learning into a standard asset pricing environment. 
Agents suspect that returns may be conditionally predictable from a large set of observable predictors. 
To exploit this perceived opportunity, they bet on sparsity and apply the apply the least absolute shrinkage and selection operator (LASSO) to select a sparse subset of variables from the high-dimensional predictor space \citep{Tibshirani1996}.

Our main result is that the very act of searching for predictability with high-dimensional statistical learning can generate short-lived pockets of return predictability. 
In equilibrium, these pockets are sparse and evolve over time, echoing recent evidence on episodic short-horizon predictability \citep{Chinco2019, Farmer2023}. 
The mechanism follows a simple life cycle. 
In finite samples, LASSO occasionally selects predictors whose estimated coefficients are large enough to survive the penalty. 
This concentrates attention on a small number of apparently strong relationships while shrinking most coefficients to zero. 
When investors trade on the resulting perceived predictability, their demand feeds back into prices, turning the perceived signal into genuine but temporary predictability. 
As new data arrive and old observations are discounted, the selected coefficients shrink back toward zero, the signal disappears, and prices revert to behavior consistent with unconditional unpredictability.

We reach this result, by formalizing this mechanism in a Lucas-type exchange economy with a single risky asset.
Under rational expectations, the price–dividend ratio is constant and returns are unpredictable. 
We follow \citet{Adam2016} and depart from rational expectations by assuming that agents know the dividend and consumption process but not the process for returns. 
Instead, the agents suspect that returns are conditionally predictable from a large set of observable signals. 
Each period, investors estimate a perceived law of motion (PLM) in which next periods return depends linearly on a sparse subset of predictors, that the agents selectes using LASSO. 
As the number of candidate predictors exceeds the sample size by a large margin, standard linear estimation techniques are infeasible. 
Beliefs then feed into equilibrium prices through the actual law of motion (ALM) governing returns.
In contrast to the linear PLM, the ALM is nonlinear and thus misspecified.
In this case, learning is unlikely to converge to the true law of motion.
Instead, the economy converges to a cross-moment consistent equilibrium (CMCE) in which agents estimated forecasting moments coincide with those generated by the equilibrium itself.
Occasional variable-selection events can then trigger predictable return episodes that persist only until the corresponding coefficients shrink back toward zero.

We test the models predictions using the cross-section of US stocks and a set of 180 news topics constrcuted in \citet{Bybee2021} using the Wall Street Journal.
Our model is able to [....]

The rest of the paper is structured as follows. 
Section 2 discusses related literature.
Section 3 presents the model environment. 
Section 4 analyzes the asset-pricing equilibrium under standard Recursive Least Squares. 
Section 5 examines Constant Gain Learning. 
Section 6 introduces LASSO learning and derives our main results regarding sparse predictability. 
Section 7 estimates the model on the cross-section of US stocks using a high-dimensional set of textual predictors. 
Section 8 concludes.

\section{Related Literature}

Our model relates to the large literature on asset pricing with learning. 
The central premise of this work is that investors have incomplete knowledge of the data-generating process governing market outcomes and therefore infer it from observed data. 
Learning is often introduced to account for empirical features, such as excess volatility and return predictability, that are difficult to replicate with fully informed agents.

Early contributions assume fully rational agents who estimate unknown parameters of a correctly specified model using least squares \citep{Marcet1989}. 
In this setting, learning converges to the rational expectations equilibrium (REE) as agents observe more data. Even in this low-dimensional environment, learning can generate excess volatility and return predictability \citep{Timmermann1993,Timmermann1996}. 
Convergence to the REE is no longer guaranteed when agents are boundedly rational; with finite-memory and finite-sample learning, beliefs may fail to converge \citep{Honkapohja2003}.

The implications become richer when agents face model uncertainty in addition to parameter uncertainty. 
\citet{barberis1998model} study investors who switch between competing forecasting rules, while \citet{hong2007simple} and \citet{Branch2010} analyze settings in which agents estimate univariate perceived laws of motion despite a multivariate true model.

More recently, \citet{adam2011internal} introduce the concept of \textit{internal rationality}, where agents optimize decisions given subjective beliefs even if those beliefs are misspecified. 
Building on this, \citet{Adam2016} embed internally rational agents in a standard asset-pricing model in which investors estimate the mapping from fundamentals to prices, generating a feedback loop that helps match key asset-pricing moments. 
Our model builds on \citet{Adam2016} by adopting their general equilibrium structure, but replaces their learning rule with a microfounded LASSO-based mechanism. 
This distinction is central: whereas their framework generates persistent belief-driven dynamics, our variable-selection mechanism delivers \textit{episodic} and \textit{sparse} predictability, consistent with evidence from high-dimensional return prediction.

Our paper also relates to research on how big data and improved statistical tools affect asset prices. 
\cite{begenau2018big} show that higher data-processing capacity improves investor forecasts and changes equilibrium prices and the cost of capital, while \cite{farboodi2020long} model investors who optimally choose whether to acquire data about fundamentals or about the behavior of other investors.

Empirical evidence indicates that return predictability is time-varying. 
Variables such as consumption growth and valuation ratios forecast returns in some periods but not in others \citep{AngBekaert2007_RFS,GoyalWelch2003_MS,LettauLudvigson2001_JF}. 
A complementary strand emphasizes sparsity: among the many predictors proposed in the literature, only a small and time-varying subset appears informative at any point in time \citep{kozak2020shrinking,freyberger2020dissecting}. 
\citet{Mclean2016} show that many return anomalies attenuate after publication, consistent with investors actively searching for and trading on these signals. 

Machine learning techniques are increasingly used to handle the ever-growing number of return predictors and to mitigate data-mining concerns \citep{Feng2020}.
\citet{Gu2020} demonstrate that machine-learning forecasts outperform traditional linear models by exploiting sparse signals and nonlinear interactions.
Most closely related to our learning mechanism is \citet{Chinco2019}, who document that LASSO uncovers short-lived pockets of high-frequency predictability concentrated in a small, shifting set of predictors. 
Finally, our work connects to initial attempts to model market outcomes when agents use statistical learning. 
\citet{Georges2021} study adaptive model selection and regularization in an agent-based market simulation. 
Relative to this simulation-based evidence, our contribution is to provide an analytical characterization of LASSO-based learning in an equilibrium environment and to derive sharp conditions under which sparse, temporary predictability arises endogenously.


% --------------------
\section{The model}
% --------------------


We study a standard \cite{Lucas1978} asset pricing model in the flavour of \cite{Adam2016}.
We consider an economy populated by a continuum of agents of mass one who trade a single risky asset.
The asset pays a stochastic dividend $D_t$ at each discrete time period $t = 0,1,2,\ldots$.
The dividend and consumption processes are driven by multivariate normally distributed shocks
\begin{equation}
    \begin{pmatrix}
    \log \varepsilon^d_t \\
    \log \varepsilon^c_t
    \end{pmatrix} \sim \mathcal{N}\left(\begin{pmatrix}
    -s^2_d/2 \\
    -s^2_c/2
    \end{pmatrix}, \begin{pmatrix}
    s^2_d & \rho s_d s_c \\
    \rho s_d s_c & s^2_c
    \end{pmatrix}\right),
\end{equation}
The dividend process follows 
\begin{equation}
    \frac{D_t}{D_{t-1}} = a \varepsilon^d_t,
\end{equation}
where $a \geq 1$ is a constant growth factor, which implies that
\begin{equation}
    \mathbb{E}_t\left[{D_{t+1}}\right] = a D_t.
\end{equation}
Similarly for the aggregate consumption process we postulate that
\begin{equation}
    \frac{C_t}{C_{t-1}} = a \varepsilon^c_t.
\end{equation}

The equilibrium price of the asset is determined by the standard no-arbitrage condition which requires that the price at time $t$ equals the discounted expected value of next period's price plus dividend, that is
\begin{equation}\label{eq:equilibrium_price}
    P_t = \delta \mathbb{E}_t \left[{ \left(\frac{C_{t+1}}{C_t}\right)^{- \gamma} P_{t+1} + D_{t+1}}\right],
\end{equation}
where $\delta$ is the discount factor and $\mathbb{E}_t[\cdot]$ is the expectation operator conditional on information available at time $t$.

Under rational expectations, the equilibrium price can be solved as
\begin{equation}
    P_t =  \frac{\delta a^{(1 - \gamma)} \rho_e}{1 - a^{(1 - \gamma)} \rho_e} D_t,
\end{equation}
where 
\begin{equation}
    \rho_e = \mathbb{E}_t \left[\left( \frac{C_{t+1}}{C_t}\right)^{- \gamma} \frac{D_{t+1}}{D_t} \right] = e ^ {\gamma(1 + \gamma) s^2_c/2} e ^ {-\gamma \rho s_c s_d} .
\end{equation}
This implies a constant price-dividend ratio and an ex-dividend gross return which is given by
\begin{equation}
    R_{t+1} = \frac{P_{t+1}}{P_t} = \frac{D_{t+1}}{D_t} = a \varepsilon^d_{t+1}.
\end{equation}
Taking logs we obtain log-returns
\begin{equation}
    r_{t+1} = \log(R_{t+1}) = \log(a) + \log(\varepsilon^d_{t+1}) \sim \mathcal{N}(\log(a) - s^2_d/2, s^2_d).
\end{equation}

Without loss of generality in what follows we set $\log(a) = s^2_d/2$. This implies that under rational expectations log-returns have mean zero and there is no time-$t$ information that can help forecasting returns.

Instead of assuming that agents know the true data generating process, we follow the adaptive learning literature and assume that agents are boundedly rational and form their expectations using econometric techniques \textcolor{red}{[citations here]}.
Equation (\ref{eq:equilibrium_price}) shows that agents have to forecast two quantities, the risk-adjusted gross return and gross dividend growth.
We follow \cite{Adam2016} in making the following assumption.

\begin{assumption}
    Agents know the process for dividends and consumption and are boundedly rational only with respect to (gross) returns.
\end{assumption}

In this way we can isolate the effects of learning about returns on asset prices.
The usual approach to learning is to endow agents with a Perceived Law of Motion (PLM) which is of the same functional form as the Actual Law of Motion (ALM) but with unknown parameters that agents estimate using available data.
In this paper we depart from the standard approach and assume that agents use linear models and try to forecast log-returns on some set of predictors available at time $t$.

% --------------------
\subsection{Univariate PLM and ALM}
% --------------------

To fix ideas we start by considering the univariate case in which agents try to predict log-returns using a single predictor $x_t$.

Specifically we assume that agents have a linear and stationary PLM for the asset return
\begin{equation}\label{eq:plm_univariate}
    r_{t+1} =  x_{t} \beta + u_t,
\end{equation}
with $x_t \sim \mathcal{N}(0, \sigma^2_x)$ being a signal available for the agent at time $t$ which they suspect might have predictive power on log-returns.

Then the agents' PLM about log-returns implies the following PLM for gross returns
\begin{equation}
    R_{t+1} = e ^ {x_{t} \beta + u_t}.
\end{equation}
This implies that gross expected returns follow a log-normal distribution and their conditional forecast is
\begin{equation}
    \mathbb{E}_t(R_{t+1}) = \mathbb{E}_t\left(\frac{P_{t+1}}{P_t}\right) = e^{x_{t} \beta + \sigma^2_u / 2} = \phi e^{x_{t} \beta},
\end{equation}
by letting $\phi = e^{\sigma^2_u / 2}$.

\begin{assumption}\label{ass:consumtion_independence}
    Within their perceived law of motion, agents believe that the return process and consumption growth are independent.
\end{assumption}

Assumption \ref{ass:consumtion_independence} allows us to derive risk-adjusted expected gross returns as
\begin{equation}
    \mathbb{E}_t \left[\left( \frac{C_{t+1}}{C_t}\right)^{- \gamma} \frac{P_{t+1}}{P_t} \right] = \mathbb{E}_t \left[\left( \frac{C_{t+1}}{C_t}\right)^{- \gamma} \right] \mathbb{E}_t\left(\frac{P_{t+1}}{P_t}\right) = a ^ {-\gamma} \phi e^{x_t \beta}.
\end{equation}
Plugging this into (\ref{eq:equilibrium_price}) yields
\begin{equation}
    P_t = \frac{\delta a^{(1 - \gamma)} \rho_e}{1 - \delta a^{- \gamma} \phi e^{x_t \beta}} D_t.
\end{equation} 
Now divide by $P_{t-1}$ to obtain actual gross returns
\begin{equation}
    R_{t} = \frac{P_{t}}{P_{t-1}} = \frac{1 - \delta a^{- \gamma} \phi e^{x_{t-1} \beta}}{1 - \delta a^{- \gamma} \phi e^{x_{t} \beta}} \frac{D_t}{D_{t-1}},
\end{equation}
and taking logs and moving one step forward we obtain the ALM 
\begin{equation}\label{eq:alm_univariate}
    r_{t+1} = \log(\varepsilon_{t+1}) + \log(1 - \kappa e^{x_{t} \beta}) - \log(1 - \kappa e^{x_{t+1} \beta}),
\end{equation}
where $\kappa := \delta a^{- \gamma} \phi$.


% --------------------
\subsection{Global well-definedness and bounded returns}
% --------------------

Expression \eqref{eq:alm_univariate} is only well-defined if
\[
1 - \kappa e^{x_t \beta} > 0
\quad \text{and} \quad
1 - \kappa e^{x_{t+1} \beta} > 0
\]
hold with probability one, which imposes state-dependent restrictions on $(x_t,\beta)$.
To avoid these domain issues and to guarantee global existence of a stochastic equilibrium, we adopt a bounded specification for the conditional mean of returns.

We modify the perceived conditional expectation of gross returns as
\begin{equation}
    \mathbb{E}\left(R_{t+1} \mid \mathcal{F}_t\right) = \phi \exp\!\Bigl( w \tanh(x_t \beta / w) \Bigr),
\end{equation}
where
\[
    w = -\log(\kappa) - \varepsilon, \qquad \varepsilon > 0.
\]
Since $\tanh(\cdot) \in (-1,1)$, this implies
\[
    w \tanh(x_t \beta / w) \in (-w,w)
    \quad \Rightarrow \quad
    e^{w \tanh(x_t \beta / w)} \in (e^{-w}, e^{w}).
\]
By construction $w < -\log(\kappa)$, so
\[
    e^{w} < e^{-\log(\kappa)} = \kappa^{-1},
\]
and therefore
\[
    1 - \kappa e^{w \tanh(x_t \beta / w)} > 0
    \quad \text{for all } x_t,\beta.
\]

The corresponding risk-adjusted expectations and price-dividend ratio are defined exactly as above, with $x_t \beta$ replaced by $w \tanh(x_t \beta / w)$.
The implied (modified) ALM is
\begin{equation}\label{eq:alm_univariate_bounded}
    r_{t+1} = \log(\varepsilon_{t+1})
    + \log\!\Bigl(1 - \kappa e^{w \tanh(x_t \beta / w)}\Bigr)
    - \log\!\Bigl(1 - \kappa e^{w \tanh(x_{t+1} \beta / w)}\Bigr).
\end{equation}

For later use, define the function
\begin{equation}\label{eq:g_beta_def}
    g_{\beta}(x) 
    := \log\!\Bigl(1 - \kappa\, e^{w \tanh(x \beta/w)}\Bigr).
\end{equation}
This function is globally well-defined and smooth in $x$ and $\beta$.

\begin{remark}
For small $\sigma_x^2$ and moderate values of $\kappa$, the original ALM \eqref{eq:alm_univariate} is well-defined with high probability.
The bounded specification \eqref{eq:alm_univariate_bounded} should then be interpreted as a globally valid extension that coincides with the baseline model in the empirically relevant region of the state space but eliminates corner cases where the log is undefined.
\end{remark}


Clearly this ALM is non-linear, which implies that the linear PLM used by agents is misspecified.
In these cases it is unlikely that the learning procedure will converge to the REE and the literature has instead focused on convergence to a Stochastic Consistent Expectations Equilibrium (SCEE) \textcolor{red}{[citations here]}.
In our case we look for an equilibirum which is cross-moment consistent in the sense that, despite their misspecification, there is a statistical self-consistency between the empirical relationship they perceive and the one which is implied by the ALM. 
Analogously to the SCEE case, we define a Cross-Moment Consistent Equilibrium (CMCE) as follows.

\begin{definition}[Cross-Moment Consistent Equilibrium]
Let $\mathcal{R}$ denote a regression operator mapping a joint distribution 
over $(r_{t+1}, x_t)$ into a parameter vector.\footnote{Examples include OLS, 
weighted least squares, LASSO, ridge, elastic net, etc.}

A pair $(\mu,\beta)$ is a \emph{Cross-Moment Consistent Equilibrium (CMCE)} 
associated with $\mathcal{R}$ if:
\begin{enumerate}
    \item $\mu$ is a non-degenerate invariant probability measure for the 
    stochastic process
    \[
        r_{t+1}
        = \log(\varepsilon_{t+1})
        + g_{\beta}(x_t)
        - g_{\beta}(x_{t+1}).
    \]
    \item The parameter $\beta$ is fixed by the regression operator applied to 
    the invariant measure, that is,
    \[
        \beta = \mathcal{R}(\mu),
    \]
    where $\mathcal{R}(\mu)$ denotes the population regression coefficient 
    obtained by applying $\mathcal{R}$ to the joint distribution of 
    $(r_{t+1}, x_t)$ induced by $\mu$.
\end{enumerate}
\end{definition}

Intuitively, while a REE imposes consistency on the \emph{entire} distributions implied by the ALM and PLM, a CMCE only requires consistency on certain moments of the two distributions, and in particular of those moments that are relevant for the agent’s forecasting problem.

We now turn to characterizing existence and stability of CMCEs under different learning algorithms.
A large empirical literature estimates return predictability models using short
rolling windows and regularized regressions such as the LASSO. In practice,
agents (or econometricians) work with limited samples, re-estimate frequently,
and penalize coefficients to avoid overfitting. This empirical behaviour implies
\emph{short memory} and time–varying parameter estimates that respond strongly
to recent observations.

In the next sections we build a theoretical benchmark that moves gradually from the
case of infinite–memory OLS, to constant–gain learning that discounts the
past, and finally to LASSO-type learning. This sequence allows us to isolate the
separate roles of (i) memory (or discounting) and (ii) sparsity–inducing regularization. 
We show that withrecursive least squares (RLS) and decreasing gain the cross–moment consistent
equilibrium (CMCE) coincides with the rational expectations equilibrium (REE)
and is locally stable under learning. When agents discount the past using a
constant gain, the estimated parameter oscillates around the CMCE/REE and never
settles. Replacing least squares by LASSO does not eliminate these fluctuations:
for small enough penalty $\lambda$ we can characterize the probability that the
LASSO estimate is different from zero, providing a rationale for temporary
``windows of predictability'' even when the underlying asset–pricing model
implies unpredictable returns. We then turn to the existence and stability of
CMCEs under the different learning mechanisms. \textcolor{red}{[citations here]}

% =====================================================
\section{Recursive Least Squares Learning}
% =====================================================

The first updating mechanism we consider is Recursive Least Squares (RLS), which is a recursive implementation of the standard Ordinary Least Squares (OLS).
In OLS, the agent estimates the coefficients of a linear regression model by minimizing the sum of squared residuals.
Assuming agents use OLS to estimate the parameters in the PLM (\ref{eq:plm_univariate}), the OLS estimator $\beta_t$ can be derived from the minimization\footnote{
We use the following timing:
at time $t$ the agent forecasts $r^e_{t+1} = x_t \beta_{t}$ which then implies a realization for $r_t$. 
Therefore, to avoid simultaneity issues, we assume that when computing the current estimate, the last available log-return observation is $r_{t-1}$.
} problem as
\begin{equation}
     \beta_t = \left(\sum^t_{i=1}  x_{i-1}^2 \right)^{-1} \sum^{t-1}_{i=1}  x_{i-1} r_i.
\end{equation}

Estimating OLS on the whole available data set is sometimes referred to as off-line estimation. 
Another approach is to estimate the model on-line, that is, as soon as a new data point arrives. 
It turns out that the estimator formula allows for a recursive formulation.
Define 
\begin{equation}
M_{t-1} = \frac{1}{t-1}  \left(\sum^{t-1}_{i=1}  x_{i-1}^2 \right), \quad \beta_t = \frac{1}{t-1} M^{-1}_{t-1} \sum^{t-1}_{i=1}  x_{i-1} r_i.
\end{equation}

First we establish a recursion for $M_t$. Start from 
\begin{equation}
\frac{t-1}{t} M_{t-1} = \frac{1}{t}  \left(\sum^{t-1}_{i=1}  x_{i-1}^2 \right) \implies \frac{t-1}{t} M_{t-1} + \frac{1}{t} x_{t-1}^2 = M_t,
\end{equation}
that is
\begin{equation}
    M_t = M_{t-1} + \frac{1}{t} \left(x_{t-1}^2 - M_{t-1} \right).
\end{equation}

For $\beta_t$ we have
\begin{equation}\label{eq:b_t_evolution}
    \begin{aligned}
        \beta_t & = \beta_{t-1} + \left[(t-1) M_{t-1} \right]^{-1} \left\{ x_{t-2} r_{t-1} + \left[(t-2) M_{t-2} \right] \beta_{t-1} \right\} - \beta_{t-1} \\
        & = \beta_{t-1} + \left[(t-1) M_{t-1} \right]^{-1} \left\{ x_{t-2} r_{t-1} - x_{t-2}^2 \beta_{t-1} \right\}.
    \end{aligned}
\end{equation}

The ALM associated with this learning mechanism is then
\begin{equation}\label{eq:alm_RLS}
    r_{t+1} = \log(\varepsilon_{t+1})
    + \log\!\Bigl(1 - \kappa e^{w \tanh(x_t \beta_t / w)}\Bigr)
    - \log\!\Bigl(1 - \kappa e^{w \tanh(x_{t+1} \beta_{t+1} / w)}\Bigr).
\end{equation}

\subsection{CMCE under RLS}
\begin{proposition}\label{prop:RLS}
    Under RLS there exist a unique CMCE given by $\beta = 0$.
    The CMCE is locally stable under learning for $\kappa \in (0,1)$.
\end{proposition}

Proof. In Appendix \ref{app:proofs}.

Notice that $\kappa = \delta a^{- \gamma} \phi$ is strictly positive for all admissible values of the parameters.
The upper bound implies a restriction on the combination of discount factor, risk aversion and dividend growth which is met for reasonable parameter values and also needed for the REE to be well-defined.



% =====================================================
\section{Weighted Least Squares and Constant Gain Learning}
% =====================================================

In Weighted Least Squares (WLS), the agent estimates the coefficients of a linear regression model by minimizing a weighted sum of squared residuals.
The weights are typically chosen to reflect the relative importance or reliability of different observations. Assuming the linear model (\ref{eq:plm_univariate}),
the WLS estimator can be derived from the minimization problem
\begin{equation}
    \min_{\beta} \sum^t_{i=1} {w_{t,i}} (y_i - \beta x_i) ^2,
\end{equation}
with solution
\begin{equation}
    {\beta} = \left(\sum^t_{i=1} {w_{t,i}} x_i^2 \right)^{-1} \sum^t_{i=1} {w_{t,i}} x_i y_i.
\end{equation}

In particular, we are interested in the case where the weights decline geometrically over time\footnote{
    In empirical applications, instead of using geometrically declining weights, it is common to use a rolling window approach, where only the most recent observations within a fixed-size window are used for estimation. 
    However, the latter is analytically intractable.
    It is possible to determine a relationship between the window size and the geometric decay parameter $\gamma$ such that the two methods give similar weights \emph{overall} on the $\ell$ most recent observations:
    $\ell \approx \log(1 - p) / \log(1 - \gamma)$, where $p$ is the proportion of total weight assigned to the $\ell$ most recent observations.
}
\begin{equation}
    w_{t,i} = \gamma (1 - \gamma)^{t-i}, \quad \gamma \in (0,1).
\end{equation}
Similar to the OLS case, a recursive relationship can be derived.

Define 
\begin{equation}
    M_t = \gamma \sum^t_{i=1} {(1 - \gamma)^{t-i}} x_i^2, \quad \beta_t = \gamma M^{-1}_t \sum^t_{i=1} {(1 - \gamma)^{t-i}} x_i y_i.
\end{equation}
First we establish a recursion for $M_t$. Start from
\begin{equation}
M_{t-1} = \gamma \sum^{t-1}_{i=1} {(1 - \gamma)^{t-1-i}} x_i^2,
\end{equation}
then
\begin{equation}
(1 - \gamma) M_{t-1} = \gamma \sum^{t-1}_{i=1} {(1 - \gamma)^{t-i}} x_i^2 \rightarrow (1 - \gamma) M_{t-1} + \gamma x_t^2 = M_t.
\end{equation}
That is
\begin{equation}
    M_t = M_{t-1} + \gamma \left(x_t^2 - M_{t-1} \right).
\end{equation}

Similarly for $b_t$ we start from
\begin{equation}
    \beta_{t-1} = \gamma M^{-1}_{t-1} \sum^{t-1}_{i=1} {(1 - \gamma)^{t-1-i}} x_i y_i,
\end{equation}
then
\begin{equation}
    (1- \gamma) \beta_{t-1} = \gamma M^{-1}_{t-1} \sum^{t-1}_{i=1} {(1 - \gamma)^{t-i}} x_i y_i ,
\end{equation}
and
\begin{equation}
 \beta_{t-1} ( M_t - \gamma x_t^2) =  \gamma \sum^{t-1}_{i=1} {(1 - \gamma)^{t-i}} x_i y_i \rightarrow \beta_{t-1} (1 - \gamma M^{-1}_t x_t^2) = \gamma M^{-1}_t \sum^{t-1}_{i=1} {(1 - \gamma)^{t-i}} x_i y_i,
\end{equation}
since \[M_{t-1} = \frac{1}{1 - \gamma} (M_t - \gamma x_t^2).\]
Adding $\gamma M^{-1}_t x_t y_t$ to both sides we obtain
\begin{equation}
    \beta_{t-1} (1 - \gamma M^{-1}_t x_t^2) + \gamma M^{-1}_t x_t y_t = \beta_t,
\end{equation}
which can be rearranged as
\begin{equation}
    \beta_t = \beta_{t-1} + \gamma M^{-1}_t x_t (y_t - x_t \beta_{t-1}).
\end{equation}

In relation to our model, this gives the same type of updating rule as RLS, with the only difference that the gain is now constant and equal to $\gamma$.
Given that the gain is not decreasing and that the updating depends on the stochastic term $\varepsilon_t$, the learning rule cannot converge to a fixed point.
However, under suitable conditions and for small values of $\gamma$, and under the same stability condition $\kappa \in (0, 1)$, it can be shown that the learning dynamics will fluctuate around the CMCE.

\begin{proposition}[Distribution under Constant Gain Learning]
    Under Weighted Least Squares learning with geometrically declining weights with parameter $\gamma$, and for $\kappa < 1$, for small values of $\gamma$ the estimated parameter $\beta$ approximately follows a normal distribution
\begin{equation}
    \beta_t \sim \mathcal{N} \left(\beta^*, \gamma s_{d}^2 \left(2 \sigma^2_x \left(1 + \frac{\kappa}{1 - \kappa} \right)\right)^{-1}\right),
\end{equation}
where $\beta^* = 0$ is the cross-moment consistent equilibrium.
\end{proposition}

% =====================================================
\section{Lasso Regression}
% =====================================================

We now consider the case in which agents use LASSO regression to estimate the parameters of their PLM.
The LASSO (Least Absolute Shrinkage and Selection Operator) estimator is defined as the solution to the optimization problem
\[
\hat{\beta}_{\text{LASSO}} = \arg\min_{\beta} \left\{\frac{1}{2}\left[(Y - X\beta)'(Y - X\beta)\right] + \lambda \sum_{j=1}^{k} |\beta_j|\right\},
\]
where $\lambda \geq 0$ is a tuning parameter that controls the amount of regularization. The LASSO estimator can be interpreted as a trade-off between fitting the data well (minimizing the sum of squared residuals) and keeping the model simple (minimizing the absolute values of the coefficients).

In general there is no closed form solution for the LASSO estimator, since the absolute value function is not differentiable at zero. However, in the special case where $X$ is an orthonormal matrix, the LASSO estimator can be expressed in closed form.
To do so we can rewrite the optimization problem as
\[
\hat{\beta}_{\text{LASSO}} 
= \arg\min_{\beta} 
\Big\{ \frac{1}{2}\left[Y'Y - 2\beta'\underbrace{X'Y}_{\beta_{\text{OLS}}} 
+ \beta'\underbrace{X'X}_{I}\beta\right]
+ \lambda \sum_{j=1}^{k} |\beta_j| \Big\},
\]
where the underbrackets follow from orthonormality and notice that the term $Y'Y$ does not depend on $\beta$, so it does not change the minimizer.
Using $A'A = \sum_{j=1}^{k} a_j^2$, we can rewrite the problem as
\[
\hat{\beta}_{\text{LASSO}} = \arg\min_{\beta} \Big\{ \sum_{j=1}^{k} \left( -\beta_j \beta_{\text{OLS},j} + \frac{1}{2}\beta_j^2 + \lambda |\beta_j| \right) \Big\}.
\]

This problem is separable in the coefficients $\beta_j$, so we can solve for each coefficient independently. The optimization problem for each $j$ coefficient is
\[\min_{\beta_j} \left\{\frac{1}{2} \beta_j^2 - \beta_j \beta_{\text{OLS},j} + \lambda |\beta_j| \right\}.\]
To solve this problem, we consider two cases:
\begin{enumerate}
    \item $\beta_{\text{OLS},j} > 0$, then $\beta_j$ has to be positive, otherwise we could decrease the objective by setting $\beta_j = 0$. 
    In this case, the problem becomes 
    \[\min_{\beta_j \geq 0} \left\{ \frac{1}{2} \beta_j^2 - \beta_j \beta_{\text{OLS},j} + \lambda \beta_j \right\},\]
    and differentiating and setting to zero gives
    \[\beta_j - \beta_{\text{OLS},j} + \lambda = 0 \Rightarrow \beta_j = \beta_{\text{OLS},j} - \lambda.\]
    Imposing non-negativity yields
    \[\hat{\beta}_{\text{LASSO},j} = \max\left(0, \beta_{\text{OLS},j} - \lambda\right).\]
    \item $\beta_{\text{OLS},j} < 0$, then $\beta_j$ has to be negative, otherwise we could decrease the objective by setting $\beta_j = 0$. In this case, the problem becomes
    \[\min_{\beta_j < 0} \left\{ \frac{1}{2} \beta_j^2 - \beta_j \beta_{\text{OLS},j} + \lambda (-\beta_j) \right\},\]    
    and
    \[\beta_j - \beta_{\text{OLS},j} - \lambda = 0 \Rightarrow \beta_j = \beta_{\text{OLS},j} + \lambda.\]
    Imposing negativity yields
    \[\hat{\beta}_{\text{LASSO},j} = \min\left(0, \beta_{\text{OLS},j} + \lambda\right).\]
\end{enumerate}

Combining both cases, we can write the LASSO estimator for each coefficient as
\begin{equation}\label{eq:lasso_from_ols}
\hat{\beta}_{\text{LASSO}} = S(\beta_{\text{OLS}}, \lambda) = \textrm{sign}(\beta_{\text{OLS}}) \cdot \max\left(0, |\beta_{\text{OLS}}| - \lambda\right),
\end{equation}
where $S(\cdot)$ is the soft-thresholding operator.

In the univariate case, we othonormality is trivially satisfied by standardizing the predictor variable, so in the discussion belowe we impose the addinitional assumption

\begin{assumption}
    The predictor variable $x_t$ is standardized to have zero mean and unit variance.
\end{assumption}
Under this condition, we can still use the recursion used in recursive and constant gain learning, and then apply the soft-thresholding operator to obtain the LASSO estimate.
Defining by $\beta^{O}$ the estimate obtained from least squares learning (either RLS or WLS), and by $\beta^{L}$ the LASSO estimate, we have
The mapping between the PLM and ALM now becomes
\begin{equation}
    \beta_{\text{OLS}} = T(\beta_{\text{LASSO}}) = - \frac{\kappa}{1 - \kappa} \beta_{\text{LASSO}},
\end{equation}
which can be solved case by case:
\begin{enumerate}
    \item If $\beta_{\text{OLS}} > \lambda > 0$, then $\beta_{\text{LASSO}} = \beta_{\text{OLS}} - \lambda \implies \beta_{\text{OLS}} =  \lambda \kappa$.
    \item If $\beta_{\text{OLS}} < -\lambda < 0$, then $\beta_{\text{LASSO}} = \beta_{\text{OLS}} + \lambda \implies \beta_{\text{OLS}} = - \lambda \kappa$.
    \item If $|\beta_{\text{OLS}}| \leq \lambda$, then $\beta_{\text{LASSO}} = 0 \implies \beta_{\text{OLS}} = 0$.
\end{enumerate}
Clearly, for the first two cases to be admissible we need $\lambda \kappa > \lambda \implies \kappa > 1$.
Given that the dynamics are still governed by the same constant gain updating rule, we can conclude that only the third case is stable in the sense that for $\kappa < 1$ the learning dynamics will fluctuate around the fixed point.
Identically to the previous case, the estimate for $\beta_{\text{OLS}}$ will be distributed as a normal distribution centered around the fixed point, with variance proportional to the gain parameter $\gamma$.
This implies that for small $\lambda$ sometimes the estimates will be large enough that the LASSO estimate is non-zero.
Given the distribution of $\beta_{\text{OLS}}$, we have:

\begin{proposition}[Probability of Non-Zero LASSO Estimate]
    Under Weighted Least Squares learning with geometrically declining weights 
    with parameter $\gamma$, and for $\kappa < 1$, for small values of $\gamma$ the probability that the LASSO estimate is different from zero is given by
\begin{equation}
    \mathbb{P}(|{\beta}_{\text{OLS}}| > \lambda) 
    = 1 - \left[\Phi\left(\frac{\lambda}{\sigma_{\text{OLS}}}\right) - \Phi\left(\frac{-\lambda}{\sigma_{\text{OLS}}}\right)\right],
\end{equation}
where $\sigma_{\text{OLS}} =  \left(\gamma s_{d}^2 \left(2 \sigma^2_x \left(1 + \frac{\kappa}{1 - \kappa} \right)\right)^{-1}\right)^{1/2}$. 
\end{proposition}


\section{Empirical Application}

The advantage of the theoretical framework developed above is that it provides a clear guide as to how to structure an empirical investigation of return predictability under learning.
We remark that when estimating our model empirically, we face the same issue that agents face in the model.
We have postulated a certain model which in theoritcal section is the actual law of motion, but in reality, given we do not know the 
true ALM, we should regard it as the modeler's perceived law of motion.
Therefore it helps to think in terms two estimations: the estimation that an agent in our model would carry out, and the estimation that we the modeler carry out.
We again endow our agents with a linear PLM
\begin{equation}
    r_{t+1} = X_t \beta_t + \eta_{t+1},
\end{equation} 

but with respect to the theoretical counterpart, we now allow for multiple predictors collected in the vector $X_t$.
While the closed-form results for the LASSO case, hold for orthonormal regressors only, in the empirical counterpart we are more lenient
and only require that predictors are uncorrelated so that equation (\ref{eq:lasso_from_ols}) does not hold deterministically but only approximately.
We then assume that agents use a rolling window, instead of geometrically declining weights, to estimate the parameters at each point in time.
Specifically at each time $t$ the regressor matrix will contain $l$ past observations, that is
\begin{equation}
    X_t = \begin{bmatrix}
    x_{t-l} \\
    x_{t-l+1} \\
    \vdots \\
    x_{t-1}
    \end{bmatrix}, \quad
    r_{t+1} = \begin{bmatrix}
    r_{t-l+1} \\
    r_{t-l+2} \\
    \vdots \\
    r_{t}
    \end{bmatrix}.
\end{equation}
For each rolling window, the agent then computes the corresponding LASSO estimator $\beta_t$.
Given this estimator, the ALM implies the modeler's PLM
\begin{equation}\label{eq:empirical_alm}
    r_{t+1} = \log(1 - \kappa e^{X_t \beta_t}) - \log(1 - \kappa e^{X_{t+1} \beta_{t+1}}) + u_{t+1}.
\end{equation}
This suggests that an emprical strategy to test for return predictability under learning is to estimate a LASSO regression on rolling windows of historical data, and then use the estimated coefficients to construct the model-implied returns.
Finally one would estimate (\ref{eq:empirical_alm}) using Non-Linear Least Squares (NLS).
However as we discussed in the theoretical section, if we as researchers estimate a linear model, we would not estimate the actual parameters used by agents, but rather
their trasformed counterpart implied by the T-map.
This suggests that the appropriate empirical strategy is, first to estimate the (mis-specified) linear model using the exact same model
we endow our agents with. Then invert the T-map to recover the parameters that agents would be using in their PLM.
Finally, use these parameters to estimate the NLS model implied by the ALM.




\bibliography{references}

\appendix
\section{Proofs}
\subsection{Proof of proposition \ref{prop:RLS}}\label{app:proofs}

In a CMCE the value of $\beta$ solves
\begin{equation}
    T(\beta) - \beta = 0,
\end{equation}
where 
\begin{equation}
    T(\beta) = \frac{\operatorname{Cov}(r_{t+1}, x_t)}{\operatorname{Var}(x_t)} = \frac{\operatorname{Cov}(r_{t+1}, x_t)}{\sigma^2_x},
\end{equation}
and $r_{t+1}$ is given by (\ref{eq:alm_univariate}) evaluated at the stationary distribution.

To compute $\operatorname{Cov}(r_{t+1}, x_t)$ we exploit the fact that $x_t$ is independent of $\log(\varepsilon_{t+1})$ and $x_{t+1}$, which implies that
\begin{equation}
    \operatorname{Cov}(r_{t+1}, x_t) = \operatorname{Cov}\left(g_{\beta}(x_t), x_t\right).
\end{equation}

By Stein's lemma
\begin{equation}
    \operatorname{Cov}(r_{t+1}, x_t) = \operatorname{Cov}(g_{\beta}(x_t), x_t) = \sigma^2_x \mathbb{E}\left[g'_{\beta}(x_t)\right].
\end{equation}
This implies that $T(\beta) = \mathbb{E}\left[g'_{\beta}(x_t)\right]$.

We can show that  $\operatorname{sign}\left[T(\beta)\right] = - \operatorname{sign}(\beta)$ unless $\beta = 0$. In fact, we have
\begin{equation}
    g_{\beta}'(x)
= -\,\frac{\kappa\,\beta\, e^{w\tanh(x \beta/w)}\, \operatorname{sech}^2(x \beta/w)}
        {1 - \kappa e^{w\tanh(x \beta/w)}}.
\end{equation}
In the numerator, every multiplicative factor except $\beta$ is strictly positive since
$\kappa>0$, $e^{w\tanh(x \beta/w)}>0$, and $\operatorname{sech}^2(x \beta/w)>0$.
The denominator is also strictly positive since $w$ is chosen accordingly.
Hence the sign of $g_\beta'(x)$ is determined entirely by $-\beta$.
In particular,
\begin{equation}
    \operatorname{sign}\left[T(\beta)\right] 
    = \operatorname{sign}\left[ - \beta \,\mathbb{E}\left(\frac{\kappa\, e^{w\tanh(x \beta/w)}\, \operatorname{sech}^2(x \beta/w)}
        {1 - \kappa e^{w\tanh(x \beta/w)}}\right)\right] = - \operatorname{sign}(\beta).
\end{equation}
Therefore $T(\beta)$ has the opposite sign of $\beta$ for all $\beta \neq 0$, and the only fixed point is $\beta = 0$.

To study local stability of the CMCE under RLS learning, we use the E-stability principle of \cite{GeorgeWEvans2001}.
To do so we show how to cast the learning dynamics in the form of a Stochastic Recursive Algorithm (SRA) of the form 
\begin{equation}\label{eq:general_sra}
    \theta_t = \theta_{t-1} + \gamma_t \mathcal{H}(\theta_{t-1}, Z_t) + \gamma_t^2 \rho_t(\theta_{t-1}, Z_t),
\end{equation}
where $\theta_t$ is the parameter vector to be estimated, $\gamma_t$ is a sequence of gains and $\mathcal{H}(\cdot)$ and $\rho_t(\cdot)$ are functions capturing the updating mechanism.
$Z_t$ is a vector of state variables that in general can be conditionally linear and evolving according to
\begin{equation}\label{eq:state_evolution_sra}
    Z_t = A(\theta_{t-1}) Z_{t-1} + B(\theta_t) W_{t},
\end{equation}
where $W_t$ is a random disturbance term. 

Under suitable assumptions about the gain sequence $\gamma_t$, the functional forms $\mathcal{H}(\cdot)$ and $\rho_t(\cdot)$, and the evolution of the state variable $Z_t$,
\cite{GeorgeWEvans2001} show that it is possible to associate to the SRA (\ref{eq:general_sra}) a deterministic Ordinary Differential Equation (ODE) of the form
\begin{equation}\label{eq:associated_ode}
    \dot{\theta} = \mathbb{E}\left[\mathcal{H}(\theta, Z)\right],
\end{equation}
where the expectation is taken with respect to the stationary distribution of $Z_t$ when $\theta$ is fixed.
The E-stability principle states that if the ODE (\ref{eq:associated_ode}) is locally stable around a fixed point $\theta^*$, then the SRA (\ref{eq:general_sra}) is also locally stable around the same fixed point.

We therefore show how to apply this framework to our model.

First we can rearrange (\ref{eq:b_t_evolution}) to get
\begin{equation}
    \begin{aligned}
        \beta_t & = \beta_{t-1} + t^{-1} M^{-1}_{t-1} \left[ x_{t-2} r_{t-1} - x_{t-2}^2 \beta_{t-1} \right] \\
        & + t^{-2} \frac{t}{t-1} M^{-1}_{t-1} \left[ x_{t-2} r_{t-1} - x_{t-2}^2 \beta_{t-2} \right].
    \end{aligned}
\end{equation}

Second, notice that $r_{t-1}$ depends on $\beta_{t-1}$ and $\beta_{t-2}$ through the ALM (\ref{eq:alm_univariate}).
We therefore define an ancillary variable $\alpha_t = \beta_{t-1}$ so that
\begin{equation}
    r_{t-1} = \log(\varepsilon_{t-1})
    + \log\!\Bigl(1 - \kappa e^{w \tanh(x_{t-2} \alpha_{t-1} / w)}\Bigr)
    - \log\!\Bigl(1 - \kappa e^{w \tanh(x_{t-1} \beta_{t-1} / w)}\Bigr).
\end{equation} 

The full updating system for the parameters is then 

\begin{equation}
    \begin{aligned}
        \beta_t = \beta_{t-1} + t^{-1} M^{-1}_{t-1} \left[ x_{t-2} r_{t-1} - x_{t-2}^2 \beta_{t-1} \right]  + t^{-2} \frac{t}{t-1} M^{-1}_{t-1} \left[ x_{t-2} r_{t-1} - x_{t-2}^2 \beta_{t-2} \right]\\
        \alpha_t = \alpha_{t-1} + t^{-1} (\beta_{t-1} - \alpha_{t-1}) + t^{-2} \left(t(t-1)\right)(\beta_{t-1} - \alpha_{t-1}) \\
        M_t = M_{t-1} + \frac{1}{t} \left(x_{t-1}^2 - M_{t-1} \right),
    \end{aligned}
\end{equation}


Now we appropriately define 
\begin{equation}
\begin{aligned}
\theta_t &= 
\begin{pmatrix}
\beta_t \\
\alpha_t \\
M_t
\end{pmatrix}, \quad
\gamma_t = t^{-1}
\end{aligned}
\end{equation}


notice that the domain of $\theta_t$ is $\mathbb{R} \times \mathbb{R} \times \mathbb{R}_{+}$, since $M_t$ is a sum of squared variables and therefore strictly positive.
We also define

\begin{equation}
\underbrace{
\begin{pmatrix}
    x_{t-1} \\
    x_{t-2} \\
    \log(\varepsilon_{t-1})
\end{pmatrix}
}_{Z_t}
=
\underbrace{
\begin{pmatrix}
    0 & 0 & 0 \\
    1 & 0 & 0 \\
    0 & 0 & 0
\end{pmatrix}
}_{A}
\underbrace{
\begin{pmatrix}
    x_{t-2} \\
    x_{t-3} \\
    \log(\varepsilon_{t-2})
\end{pmatrix}
}_{Z_{t-1}}
\;+\;
\underbrace{
\begin{pmatrix}
    1 & 0 \\
    0 & 0 \\
    0 & 1
\end{pmatrix}
}_{B}
\underbrace{
\begin{pmatrix}
    x_{t-1} \\
    \log(\varepsilon_{t-1})
\end{pmatrix}
}_{W_t}.
\end{equation}

The state evolution clearly satisfies assumpotions (B.1) and (B.2) in Section 6.2.1 of \cite{GeorgeWEvans2001},
since both $x_t$ and $\log(\varepsilon_t)$ are normal i.i.d variables with finite moments of all orders.
Assumption (A.1) is satisfied by construction of the gain sequence $\gamma_t = t^{-1}$.
For assumotion (A.2) and (A.3) write 

\begin{equation}
    \begin{aligned}
        \mathcal{H}(\theta_{t-1}, X_t) = \begin{pmatrix}
        M^{-1}_{t-1} \left[ x_{t-2} r_{t-1} - x_{t-2}^2 b_{t-1} \right] \\
        b_{t-1} - a_{t-1} \\
        x_{t-1}^2 - M_{t-1}
    \end{pmatrix}, \\
    \end{aligned}
\end{equation}

\begin{equation}
    \begin{aligned}
        \rho_t(\theta_{t-1}, X_t) = \begin{pmatrix}
        \frac{t}{t-1} M^{-1}_{t-1} \left[ x_{t-2} r_{t-1} - x_{t-2}^2 b_{t-2} \right] \\
        (t(t-1))(b_{t-1} - a_{t-1}) \\
        0
    \end{pmatrix}. \\
    \end{aligned}
\end{equation}


To verify Assumption (A.2), we observe that each component of $\mathcal{H}(\theta_{t-1},Z_t)$ is continuous in $\theta$ on any compact subset of the parameter space. The transformation $w\tanh(x\theta/w)$ is smooth and bounded, which implies that
\[
\log\!\left(1-\kappa e^{\,w\tanh(x\theta/w)}\right)
\]
remains uniformly bounded for all $(\theta,x)$ whenever $\kappa e^{w}<1$. Consequently, the return innovation $r_{t-1}$ is bounded by a constant plus $|\log\varepsilon_{t-1}|$, and the expressions $x_{t-2}r_{t-1}$, $x_{t-2}^2 b_{t-1}$, and $x_{t-1}^2$ grow at most polynomially in $|x_t|$. Since the exogenous process $(x_t)$ is strictly stationary and ergodic with $\mathbb{E}[x_t^2]<\infty$ and $\mathbb{E}[x_t^4]<\infty$, all expectations defining $h(\theta)=\mathbb{E}[\mathcal{H}(\theta,X_t)]$ exist and $\mathcal{H}$ satisfies the polynomial-growth and continuity requirements of Assumption (A.2).

Assumption (A.3) holds provided that $\mathcal{H}(\theta,Z)$ is twice continuously differentiable in $\theta$ with second derivatives bounded on every compact set $Q$. In our case, the mappings
\[
\theta \mapsto w\tanh(x\theta/w), \qquad 
\theta \mapsto \log\!\left(1-\kappa e^{\,w\tanh(x\theta/w)}\right)
\]
are $C^\infty$ with all derivatives bounded uniformly on compacts, because $\tanh(\cdot)$ and its derivatives are bounded and the restriction $\kappa e^{w}<1$ keeps the logarithmic term uniformly away from its singularity. It follows that each component of $\mathcal{H}$ is twice continuously differentiable in $\theta$ with bounded second derivatives on compact sets, which verifies Assumption (A.3).

The E-stability principle then applies and we can study the stability of the equilibrium by analyzing the associated ODE.
We first obtain 

\begin{equation}
    \begin{aligned}
        h(\theta) = \lim_{t \to \infty} \mathbb{E} \left[\mathcal{H}(\theta, \Bar{Z}_t)\right] = \begin{pmatrix}
        M^{-1} \left[ \mathbb{E}(x_{t-2} r_{t-1}) - \beta \sigma^2_x \right] \\
        \beta - \alpha \\
        \sigma^2_x - M
    \end{pmatrix} =
    \begin{pmatrix}
        M^{-1}\left[(T(\alpha) - \beta) \sigma^2_x\right] \\
        \beta - \alpha \\
        \sigma^2_x - M
    \end{pmatrix}
    \\
    \end{aligned}
\end{equation}
where $\Bar{Z}_t$ refers to the evolution of the state variable, under fixed values of the parameters.

Then we evaluate the stability of the cross-moment consistent equilibrium by studying the associated ordinary differential equation (ODE)
\begin{equation}
    \frac{d \theta}{d \tau} = h(\theta).
\end{equation}
Therefore we compute the Jacobian matrix of $h(\theta)$ evaluated at the equilibrium point $\theta = (0,0,\sigma^2_x)$

\begin{equation}
    J = \begin{pmatrix}
         - 1 & T'(0) & 0 \\
        1 & -1 & 0 \\
        0 & 0 & -1
    \end{pmatrix}.
\end{equation}

The eigenvalues of $J$ are given by
\begin{equation}
    \lambda_{1,2} = -1 \pm i \sqrt{\frac{\kappa}{1 - \kappa}}, \quad \lambda_3 = - 1,
\end{equation}

since 
\begin{equation}
    T'(\alpha) = -\mathbb{E}\left[\frac{\kappa x^2 e ^ {w \tanh(x \alpha / w)} \textrm{sech}^2(x \alpha / w)}{1 - \kappa e ^ {w \tanh(x \alpha / w)} }  \right] \left(\sigma^2_x\right)^{-1} .
\end{equation}

therefore the ODE is asymptotically stable for the whole range of $\kappa \in (0,1)$ . 


\end{document}